\section*{Bab \ref{ch:b3:organometallic-bonding} --- Kimia Organologam: Ikatan dan Ligan}

\begin{solution}{exo:b3:organometallic-bonding:1}
$\ce{Fe(CO)5}$: $8 + 10 = 18$. $\ce{Mn2(CO)10}$: $7 + 10 + 1 = 18$ per Mn. $\ce{CpMo(CO)3H}$: $6 + 5 + 6 + 1 =
18$. $\ce{[PtCl3(C2H4)]-}$: $10 + 3 + 2 + 1 = 16$. $\ce{Cr($\eta^6$-C6H6)(CO)3}$: $6 + 6 + 6 = 18$. $\ce{Cp2ZrCl2}$:
$4 + 10 + 2 = 16$.
\end{solution}

\begin{solution}{exo:b3:organometallic-bonding:2}
Fe(0) $d^8$; Mn(0) $d^7$; Mo(II) $d^4$; Pt(II) $d^8$; Cr(0) $d^6$; Zr(IV) $d^0$.
\end{solution}

\begin{solution}{exo:b3:organometallic-bonding:3}
$\ce{[Mn(CO)6]+} > \ce{Cr(CO)6} > \ce{[V(CO)6]-}$: makin kaya elektron logamnya, makin banyak ia
berdonasi balik ke $\pi^*$ dan makin lemah ikatan-ikatan C--O-nya.
\end{solution}

\begin{solution}{exo:b3:organometallic-bonding:4}
Etana ($\ce{Mn(CO)5}$ dan $\ce{CpFe(CO)2}$ \omterm{def:b3:organometallic-bonding:isolobal}{isolobal} dengan \ce{CH3}); etana lagi; tetrahedrana
$\ce{(CH)4}$ yang tiga CH-nya diganti dengan $\ce{Co(CO)3}$.
\end{solution}

\begin{solution}{exo:b3:organometallic-bonding:5}
$fac$ ($C_{3v}$): dua (A$_1$ + E). $mer$ ($C_{2v}$): tiga (2A$_1$ + B$_1$).
\end{solution}

\begin{solution}{exo:b3:organometallic-bonding:6}
Fosfina yang meruah menyesaki logam: disosiasi meringankan tegangan itu, sehingga
berlangsung lebih cepat, dan zat antara berkoordinasi rendah lebih disukai.
\end{solution}

\begin{solution}{exo:b3:organometallic-bonding:7}
Kompleks kromium itu adalah \omterm{def:b3:organometallic-bonding:carbene}{karbena Fischer} (substituen OMe, Cr(0), ligan CO): sebuah
amina menyerang karbon \omterm{def:b3:organometallic-bonding:carbene}{karbena} yang elektrofilik dan menggantikan OMe. Alkilidena
tantalum itu adalah \omterm{def:b3:organometallic-bonding:carbene}{karbena Schrock}: karbonnya yang nukleofilik menyerang karbonil sebuah
keton, sehingga terbentuk alkena dan satuan Ta=O, seperti pada reaksi Wittig.
\end{solution}

\begin{solution}{exo:b3:organometallic-bonding:8}
Konformasi goyang: tidak ada tumpang tindih $\delta$, kedua elektron $\delta$ menjadi nonikatan: orde ikatan 3. $\ce{[Re2Cl8]^{4-}}$:
$\sigma^2\pi^4\delta^2\delta^{\ast2}$, orde ikatan 3.
\end{solution}

\begin{solution}{exo:b3:organometallic-bonding:9}
Jarak M$\cdots$H yang pendek (difraksi, terbaik difraksi neutron); bilangan gelombang uluran C--H
yang rendah; sinyal NMR pada medan tinggi dengan kopling C--H satu ikatan yang berkurang.
\end{solution}

\begin{solution}{exo:b3:organometallic-bonding:10}
16: $d^8$ segi empat datar ($\ce{[PtCl4]^{2-}}$, katalis Wilkinson), yang orbital untuk elektron ke-18-nya
terlalu tinggi; \omterm{def:b3:organometallic-bonding:metallocene}{metalosena} logam awal seperti $\ce{Cp2ZrCl2}$, yang terlalu sesak untuk menampung ligan
tambahan. 17: $\ce{V(CO)6}$, yang tidak dapat berdimerisasi karena alasan sterik. 19: kobaltosena,
yang elektron tambahannya menempati orbital antiikatan (dan mudah lepas).
\end{solution}

\begin{solution}{exo:b3:organometallic-bonding:11}
CO adalah akseptor $\pi$ yang kuat: CO menurunkan tingkat $t_{2g}$ dan membuat $\Delta_{\mathrm o}$ besar, jauh di atas
energi pemasangan: $t_{2g}^6$, diamagnetik. Transisi $d$--$d$, pada $\Delta_{\mathrm o}$, jatuh di daerah ultraviolet,
dan tidak ada cahaya tampak yang diserap.
\end{solution}

\begin{solution}{exo:b3:organometallic-bonding:12}
Donasi dari $\pi$ dan donasi balik ke $\pi^*$ sama-sama menurunkan orde ikatan C--C. Dalam garam
Zeise (Pt(II), yang donasi baliknya sedang), ikatannya sedikit memanjang; logam bervalensi rendah
yang kaya elektron berdonasi balik dengan kuat dan kompleksnya mendekati limit
\omterm{def:b3:organometallic-bonding:dcd}{metalasiklopropana}, dengan ikatan C--C yang mendekati ikatan tunggal.
\end{solution}

\begin{solution}{pb:b3:organometallic-bonding:1}
\textbf{1.} $\ce{Fe^{2+}}$ ($d^6$) $+ 2 \times 6 = 18$.
\textbf{2.} $8 + 2 \times 5 = 18$.
\textbf{3.} $9 + 10 = 19$.
\textbf{4.} $10 + 10 = 20$.
\textbf{5.} $18 - 1 = 17$.
\textbf{6.} Fe(II) $d^6$; Co(II) $d^7$; Ni(II) $d^8$; Fe(III) $d^5$.
\textbf{7.} $a_{1g}$ ($d_{z^2}$), $e_{1g}$ ($d_{xz}$, $d_{yz}$), $e_{2g}$ ($d_{xy}$, $d_{x^2-y^2}$).
\textbf{8.} Pasangan $e_{1g}$ bertumpang tindih kuat dengan kombinasi cincin $e_{1g}$ yang terisi: orbital
ikatan yang sebagian besar berada pada cincin, dan orbital antiikatan $e_{1g}^\ast$ yang sebagian besar berada pada logam, yang tinggi
energinya.
\textbf{9.} $e_{2g} \approx a_{1g} < e_{1g}^\ast$.
\textbf{10.} $e_{2g}^4a_{1g}^2$: tidak ada elektron yang tidak berpasangan.
\textbf{11.} Satu elektron di $e_{1g}^\ast$: satu elektron tidak berpasangan, $1.73\,\mu_B$.
\textbf{12.} Dua elektron di kedua orbital $e_{1g}^\ast$, dengan spin sejajar: dua elektron tidak berpasangan, $\sqrt{8} =
2.83\,\mu_B$.
\textbf{13.} $e_{2g}^3a_{1g}^2$: satu elektron tidak berpasangan. Lubang $e_{2g}^3$ memberikan keadaan yang
terdegenerasi secara orbital (E) dengan \omterm{def:b3:complex-spectra-magnetism:moment}{kontribusi orbital}.
\textbf{14.} Elektron ke-19-nya berada di orbital antiikatan dan mudah dilepaskan,
sehingga terbentuk kation kobaltosenium 18 elektron.
\textbf{15.} Elektron yang dilepaskan berasal dari orbital yang hampir nonikatan:
strukturnya hampir tidak berubah (\omterm{def:b3:complex-mechanisms:reorganisation}{energi reorganisasi} kecil, \cref{ch:b3:complex-mechanisms}).
\textbf{16.} Pasangan itu cepat dan reversibel dalam kebanyakan pelarut, senyawanya
larut dan stabil, dan potensialnya hanya sedikit bergantung pada pelarut karena
muatannya tersebar pada ion yang besar.
\textbf{17.} Dengan 20 elektron, dua di antaranya antiikatan, nikelosena bereaksi dengan melepaskan kedua elektron itu:
senyawa itu mengikat ligan disertai pergeseran atau lepasnya sebuah cincin, atau teroksidasi.
\textbf{18.} $\ce{CpFe(CO)2}$: $8 + 5 + 4 = 17$ elektron, kurang satu dari 18, satu orbital batas: \omterm{def:b3:organometallic-bonding:isolobal}{isolobal}
dengan \ce{CH3}. Dimer $\ce{[CpFe(CO)2]2}$ adalah ``etana'', dengan ikatan Fe--Fe (dalam praktik, dua ligan CO
menjembataninya).
\textbf{19.} $\ce{CpFe(CO)2-Mn(CO)5}$, dengan ikatan Fe--Mn.
\textbf{20.} $\ce{PPh3}$ adalah donor yang lebih baik dan akseptor $\pi$ yang lebih buruk daripada CO: logam berdonasi balik
lebih banyak ke CO yang tersisa, yang bilangan gelombang ulurannya turun.
\textbf{21.} Dua (A$_1$ + E).
\textbf{22.} Pergeseran cincin ke $\eta^3$ membebaskan koordinasi senilai dua elektron: ligan yang masuk
dapat terikat secara asosiatif tanpa logam melampaui 18 elektron.
\textbf{23.} \textbf{$\mu_{\text{spin saja}}(\text{nikelosena}) = \sqrt{2 \times 4} = 2.83\,\mu_B$.}
\end{solution}
